\documentclass{article}
\usepackage[margin=25mm]{geometry}
\usepackage{amsmath,amssymb}
\usepackage{xurl}
\usepackage[colorlinks=true,linkcolor=blue,citecolor=blue,urlcolor=blue]{hyperref}
% Shared commands for notes in docs/paper-gaps/.

\DeclareUnicodeCharacter{2080}{\ifmmode{}_{0}\else\textsubscript{0}\fi}
\DeclareUnicodeCharacter{2081}{\ifmmode{}_{1}\else\textsubscript{1}\fi}
\DeclareUnicodeCharacter{2082}{\ifmmode{}_{2}\else\textsubscript{2}\fi}
\DeclareUnicodeCharacter{2099}{\ifmmode{}_{n}\else\textsubscript{n}\fi}

\newcommand{\Lean}{\textsc{Lean}}
\ExplSyntaxOn
\NewDocumentCommand{\leanid}{m}{
  \ifmmode
    \textsf{\detokenize{#1}}
  \else
    \group_begin:
    \urlstyle{sf}
    \tl_set:Nn \l_tmpa_tl {#1}
    \tl_replace_all:Nnn \l_tmpa_tl {\_} {_}
    \exp_args:NV \path \l_tmpa_tl
    \group_end:
  \fi
}
\ExplSyntaxOff
\providecommand{\tr}{\operatorname{tr}}
\newcommand{\tnleanIssueBase}{https://github.com/LionSR/TNLean/issues/}
\newcommand{\tnleanPrBase}{https://github.com/LionSR/TNLean/pull/}
\newcommand{\tnleanDocBase}{https://sirui-lu.com/TNLean/docs/}
\newcommand{\ghissue}[1]{\href{\tnleanIssueBase#1}{GitHub issue \##1}}
\newcommand{\ghpr}[1]{\href{\tnleanPrBase#1}{PR \##1}}
\newcommand{\doclink}[1]{\href{\tnleanDocBase#1.html}{\path{#1}}}

% Dirac notation and the identity operator, as in blueprint/src/macros/common.tex.
% \providecommand keeps note-local definitions authoritative.
\usepackage{dsfont}
\providecommand{\ket}[1]{|#1\rangle}
\providecommand{\bra}[1]{\langle #1|}
\providecommand{\braket}[2]{\langle #1 | #2 \rangle}
\providecommand{\ketbra}[2]{|#1\rangle\!\langle #2|}
\providecommand{\Id}{\mathds{1}}
\providecommand{\one}{\mathds{1}}
\providecommand{\C}{\mathbb{C}}
\providecommand{\MN}[1]{M_{#1}(\C)}
\providecommand{\MD}{M_D(\C)}

% Verdict marker: \gapnote{<kind>}{<status>}. Kind is the policy
% classification (clarification, local-correction, scope-restriction,
% unfaithful, false-source, open-gap); status is open, wip (elimination
% underway), resolved, or historical. Machine-read by the site index;
% prints a verdict line.
\newcommand{\gapnote}[2]{\par\noindent\textbf{Verdict:} #1 (#2).\par}

\providecommand{\gapnote}[2]{\par\noindent\textbf{Verdict:} #1 (#2).\par}
\title{Independent Parameters in the ALVV Ground-State Approximation}
\author{}
\date{November 3, 2026}
\begin{document}
\maketitle
\gapnote{local-correction}{wip}

\section{The quantitative conclusion}

Theorem~4.9 of Arad, Landau, Vazirani and Vidick~\cite{Arad2016RigorousRG},
\href{https://arxiv.org/abs/1602.08828v2}{arXiv:1602.08828v2}, asserts a
fixed-cut Schmidt-rank bound of the form
\begin{align}
    s(\delta)\le r\exp\!\left[\widetilde O\!\left(
        \gamma^{-1}\log^3d+
        \gamma^{-1/4}\log^{3/4}(1/\delta)\log d\right)\right].
    \label{eq:alvv_approximation_rank}
\end{align}
The subsequent global matrix product state bound replaces $\delta$ by
$\delta/n$. This note concerns the parameter choice in that proof, not a
counterexample to either conclusion.
The source is \path{local_alg_arxiv_v2.tex}, line 1080; the relevant
independent-parameter estimate is Corollary~4.6,
\path{c:goodagsphardtrunc}, lines 979--1005.

\section{The source assumptions and the normalized gap}

Assumption (DG) in the source, lines 275--279, requires a ground energy
$\epsilon_0$ and a second smallest distinct eigenvalue $\epsilon_1$, with
$\epsilon_1-\epsilon_0\ge\gamma$. Thus the distinct spectral values in the
formal statements express a source hypothesis. Scalar Hamiltonians are
outside this stated setting.

The local normalization $0\le h_i\le\Id$ does not force the gap to be at
most one. On three qubits, put
$P=\Id\otimes\ketbra{1}{1}\otimes\Id$ and take both nearest-neighbour
terms to equal $P$. The first is supported on sites $1,2$ and the second
on sites $2,3$. Each term is a positive contraction, while
$H=2P$ has spectrum $\{0,2\}$, ground-space dimension four and gap two.
This corrects the normalization comparison in tracker~\ghissue{9809}.
ALVV states the local normalization at source lines 267 and 806; no
global assumption $\gamma\le1$ is stated there or in Assumption (DG).

The bounds below use $g=\min\{\gamma,1\}$. Decreasing a supplied gap
lower bound to $g$ is always legitimate. Comparing the resulting bound
with the printed inverse powers of the original $\gamma>1$ requires an
additional uniform spectral bound or a separate large-gap argument.
That comparison remains open here.

\section{The printed choice is not uniform in the gap}

Corollary~4.6 gives the decay exponent
\begin{align}
    E=\frac{\ell^2\sqrt\gamma}{\sqrt{t+\ell}}.
    \label{eq:alvv_approximation_exponent}
\end{align}
The proof of Theorem~4.9 prescribes
$t=\Theta(\gamma^{-1/4}\log(1/\delta))$ and
$\ell=\Theta(t^{3/4})$.
If the constants in these comparisons are universal, this does not make
$E$ a sufficiently large multiple of $\log(1/\delta)$ for all gaps.
For example, let $g=\gamma\downarrow0$ and
$\delta=\exp(-g^{-2})$. Then
\begin{align}
    t      & =\Theta(g^{-9/4}),                                &
    \ell   & =\Theta(g^{-27/16}),                              &
    \ell/t & \longrightarrow0,\notag                             \\
    E      & =\Theta(g^{-7/4})=o\!\left(\log(1/\delta)\right).
    \label{eq:alvv_approximation_prescription}
\end{align}
This regime is compatible with inverse-polynomial accuracy by taking
$n$ of order $\exp(g^{-2})$. Rounding $n$ or $\ell$ changes none of these
comparisons. Thus the corollary's decay estimate does not justify the
claimed accuracy from this prescription with universal constants.
The obstruction concerns the prescription; it does not disprove
(\ref{eq:alvv_approximation_rank}).

\section{A choice retaining the two scales}

Normalize the gap to $0<g\le1$. For a target scale $0<\tau\le1$, put
$u=1+\log(1/g)+\log(1/\tau)$. There is a universal constant $C\ge21$
for which $t=Cu$ satisfies
\begin{align}
    288e^{11/8}e^{-(t-g)/8}\le\frac{\tau g}{10^6}.
    \label{eq:alvv_approximation_tail}
\end{align}
The actual centered five-region Hamiltonian, with three bulks clipped at
$t$ and windows of length $2\ell$, has norm at most $3t+4\ell$.
For any $a\ge0$ with $a\le t$, choose
\begin{align}
    b    & =(7ta^2/g)^{1/4},              &
    \ell & =\left\lceil7/g+b\right\rceil.
    \label{eq:alvv_approximation_choice}
\end{align}
Then $\ell\ge1$, $7\le\ell g$, $7ta^2\le g\ell^4$, and
$\ell\le8/g+b$. These inequalities control both regimes in the
physical norm budget. For $0<R\le3t+4\ell$:
if $\ell\le t$, then $R\le7t$ and $a^2R\le g\ell^4$;
if $t\le\ell$, then $a\le\ell$, $R\le7\ell$ and
$a^2R\le7\ell^3\le g\ell^4$.
Consequently
\begin{align}
    a\le\ell^2\sqrt{g/R}.
    \label{eq:alvv_approximation_corrected_exponent}
\end{align}
In particular, when $t=Ca$, the same integer choice satisfies
\begin{align}
    \ell\le\frac8g+(7C)^{1/4}g^{-1/4}a^{3/4}.
    \label{eq:alvv_approximation_quarter_power}
\end{align}
The low-cluster spectral separation is $4g/5$ rather than $g$.
Using $4g/5$ in (\ref{eq:alvv_approximation_choice}) only changes the
universal constants: the inverse-gap term becomes $10/g$ and the
fourth-root coefficient becomes $(35C/4)^{1/4}$.
Thus both logarithms in $u$ are retained, while the accuracy dependence
has the power required in (\ref{eq:alvv_approximation_rank}).

\section{The resulting fixed-cut estimate}

Write $r=\dim G$, $u=1+\log q+\log(1/g)$ and
$a=16+\log(1/g)+\log(1/\delta)$.
The normalized physical approximation now gives one universal constant $K>0$
such that every original unit $\psi\in G$ has a unit approximant $\phi$ with
$|\langle\phi,\psi\rangle|\ge1-\delta$ and
\begin{align}
    \log\operatorname{SR}_{J:J^c}(\phi)
    \le\log r+K\bigl(g^{-1}u^3+g^{-1/4}a^{3/4}u\bigr)(1+\log a).
    \label{eq:alvv_approximation_explicit_rank}
\end{align}
The initial viable subspace is fixed before the accuracy and vector are
chosen. Its dimension contributes $rD_0^{24}$, and the independent
polynomial family contributes $D_1^2$. Thus the integer rank budget is
$rD_0^{24}D_1^2$, retaining only one factor $r$.
The vector is normalized after taking the viable-space projection witness
and is transported to the original Hilbert space by the actual cut isometry.
No rank or overlap assumption is supplied.
The construction covers intervals, including empty intervals; a prefix
interval gives the source's fixed cut.

The gap in (\ref{eq:alvv_approximation_explicit_rank}) is explicitly
$g=\min(\gamma,1)$. For $0<\gamma\le1$, the estimate has the source's
quarter-power gap and three-quarter-power accuracy dependence, with the
logarithmic correction displayed. For larger supplied gaps it uses $g=1$;
it does not assert an improvement proportional to $\gamma^{-1}$ or
$\gamma^{-1/4}$ in that regime. Comparing the capped estimate with the
unrestricted printed quantitative expression requires a universal upper
bound on a nontrivial nearest-neighbour ground-space gap, or a separate
large-gap argument. That comparison is not proved here.
The closed spectral cutoff retaining the entire original ground space is
as documented in the companion
\href{https://sirui-lu.com/QICLean/paper-gaps/alvv16_spectral_projector_boundary.pdf}
{spectral-projector note}.

\section{One global matrix product state}

Put $b=16+\log(1/g)+\log(16(n+1)/\delta)$.
The simultaneous-cut construction chooses an integer $B\ge1$ before
any ground vector, with one universal $K>0$ such that
\begin{align}
    \log B \le K\bigl(g^{-1}u^3+g^{-1/4}b^{3/4}u\bigr)(1+\log b).
    \label{eq:alvv_approximation_global_bond}
\end{align}
Every original unit $\psi\in G$ has a unit open-boundary matrix product
approximant $\Phi$ of bond dimension at most $rB$, absolute overlap at least
$1-\delta$, and squared norm error at most $\delta/2$.

The final vector is not assembled from the separate fixed-cut approximants.
Those vectors are used only to bound the original marginal's discarded mass
by $2\delta_c$, where $\delta_c=\delta/(16(n+1))$.
The single product of the original marginals' leading spectral projections
has squared norm error at most the sum of those discarded masses, hence
at most $\delta/8$. Its normalization gives $\Phi$ and preserves every
bond bound. Choosing $B=\lceil e^E\rceil$ for the common logarithmic rank
budget costs at most one in $\log B$ and retains the single factor $r$.

The remaining quantitative comparison concerns the cap $g=\min(\gamma,1)$
for larger supplied gaps, as discussed above. The low-energy extension and
the binary-tree algorithm are separate results and are not asserted here.


\section*{The bounded recursive witness construction}

For the same original full ground space, a balanced recursion can be
constructed with full local spaces at leaves of length at most $4\ell$.
At every longer parent, one actual family is chosen from the Hamiltonian
before the sampling size and the child spaces. Its literal ordered words
act on the actual Haar coordinate sample, followed by one collective trim.
Writing $D_{\rm ref}=r^2B$, take
\[
  A=1+\log r+\log D+\log(1/\beta),\qquad
  s=\lceil4096rA\rceil,\qquad
  M=D^6s^3,\qquad
  \xi=(1000D_{\rm ref}\sqrt{n+1})^{-1}.
\]
The corrected fixed-input failure bound is at most $\beta<1$, uniformly
for parent dimensions at most $D^{12}s^6$. Choosing successful unitaries
therefore gives actual recursively selected spaces, all $3/200$-viable
for the original ground space. Each has at most $M$ recorded generators,
whose internal-cut Schmidt ranks are at most
$10^6D_{\rm ref}^2(n+1)M$. This bound concerns the recorded generators,
not arbitrary linear combinations in their span. The whole exterior is
retained; the squared reference factor is not replaced by $D_{\rm ref}$.

The independent scalar tail and filter conditions remain hypotheses of
this bounded result. It uses block leaves and one final trim, whereas
Figure~\texttt{fig:merge-prime} of the source places trimming within the
error-reduction rounds. It proves existence of one recursively selected
tree, not a joint probability theorem for the adaptive algorithm. A
pointwise choice of a measurable hull of each bad set does not by itself
supply measurable dependence on the earlier spaces. The source theorem
\texttt{t:alg1} remains separate: its full parameter selection, modified
merge analysis, adaptive probability, energy-estimation invariant and
computational cost remain separate. The final spectral extraction for
the modified witness is established in the refinement below.


\section*{Automatic parameters for this witness}

The numerical tail and filter conditions in the bounded witness above can
be derived uniformly from the physical input. Choose a universal $C\ge1$
from the scalar parameter theorem, before the chain and Hamiltonian, and put
\[
  m=\max\{q,50\},\quad v=1+\log m+\log(1/g),\quad
  \ell=\lceil Cv^2/g\rceil,\quad D=(m\ell^2)^{22\ell}.
\]
Then $\ell\ge21$ and
\[
  \log D\le22(C+1)(6+2\log(C+1))v^3/g.
\]
Take $t=\ell$ and $\Delta=10^{-5}D^{-12}$. The scalar contraction
estimate holds uniformly up to $8\ell$, hence throughout the required
range $R\le3t+4\ell=7\ell$. The truncation tail estimate is at most
$g/(10^6m)$, which supplies the required extra factor $1/m$.
Thus the physical recursive witness has a version with no supplied
numerical tail or filter hypotheses. The families still precede every
sampling size and child space.

This parameter substitution does not identify the witness with the
printed modified merge procedure. Its block leaves, final trim, capped
gap and ranks of recorded generators remain as stated above. It does
not by itself prove an adaptive joint success probability, the
energy-estimation invariant, final spectral extraction or the claimed
running time. The separate refinement below supplies the spectral
extraction for a common-threshold version of the modified tree.


\section*{Native matrix product states at the actual root}

The actual recursively selected root space can be read on the original
chain coordinates. Its regional internal cuts are the native chain cuts.
If its at most $M$ recorded generators have ranks at most
$R=10^6D_{\rm ref}^2(n+1)M$, the finite-span estimate gives bond at most
\[
  MR=10^6D_{\rm ref}^2(n+1)M^2
\]
for every vector in this actual root space. The extra factor $M$ is
retained: the per-generator bound is not assigned unchanged to the whole
span. The families, actual recursive operations and original-ground-space
viability at every intermediate interval are the same as above. No
independent matrix product approximation is substituted for the root.

This constant-accuracy conclusion alone does not perform final spectral
extraction or establish running time or adaptive joint success probability.
The normalized witness below uses whole-region viability of this actual
root, with its coordinate and norm identifications proved.

\section*{Normalized witnesses in this same root}

The full interval has an empty exterior. Its viability inequality is
therefore exactly the closeness inequality for the native image $S$ of
the same recursively constructed root. Let $P_S$ be its orthogonal
projection. For every unit $\Gamma\in G$, this gives
$\|P_S\Gamma\|^2\ge197/200$. The actual normalized projection
\[
  \phi=\frac{P_S\Gamma}{\|P_S\Gamma\|}\in S
\]
has squared overlap at least $197/200$ and squared distance at most
$3/100$, with $\langle\phi,\Gamma\rangle=\|P_S\Gamma\|$ real and
nonnegative. It retains the all-span native MPS bound
$10^6D_{\rm ref}^2(n+1)M^2$.
The operator families and root precede the choice of $\Gamma$.
This is an existential constant-accuracy witness in that root. Constant
viability alone does not give an accurate ground-space basis without
knowing $\Gamma$. The following construction supplies the required
refinement and extraction for the modified common-threshold tree.


\section*{Refinement and Ritz extraction at one common threshold}

Let $\mu$ be the original ground energy, $\gamma$ the gap, and $r$ the
actual ground-space dimension. The same physical improving families can
be selected before every requested accuracy $\eta>0$. Put
\[
  A=\|H\|-\mu,\qquad c=\gamma/A,\qquad
  \rho=\frac{\eta\gamma}{2A+\eta\gamma},\qquad
  \varepsilon=\frac{c^2\rho}{32768}.
\]
Choose the physical reference space at accuracy $\varepsilon$, with
all-span MPS bond bound $b$, and let $b_0$ be the reference bound used in
the tree. With $D$ the common filter count parameter, choose
\[
  b_* =\max\{b_0,b\},\qquad
  N=\left\lceil\frac{100D^{12}}{c\rho}\right\rceil,\qquad
  \xi=\frac1{1000b_*N\sqrt{n+1}}.
\]
This threshold satisfies both the tree bound and the final refinement
budgets, including the strict inequality
$\sqrt{nrb}\,\xi<\eta/(10000D^{12})$.
Fix $0<\beta<1$, and put
\[
  s=\left\lceil4096r\bigl(1+\log r+\log D+\log(1/\beta)\bigr)\right\rceil,
  \qquad M=D^6s^3.
\]
Construct the actual balanced tree using this same $\xi$ at every long
parent. Write $S_0$ for its root; $M$ bounds the number of its recorded
generators.

Starting from this actual root, apply the literal raw filter
$K=I-H/\|H\|$ followed by collective trimming at $\xi$, for
\[
   j_* =\left\lceil(4/c)\log(1/\rho)\right\rceil+1
\]
steps. The retained space $W$ is $\rho$-close to the original full ground
space, and its dimension is at most $M$. Let $R$ be the actual lowest
$r$-dimensional Ritz space of the compression of the original $H$ to
$W$. Then $R\subseteq W$, $\dim R=r$, and its two projection errors
relative to the original ground space are at most $\eta/2$.
Indeed the projected trial space has excess energy at most
\[
   \frac{A\rho}{1-\rho}=\frac{\eta\gamma}{2},
\]
and the original gap and the variational principle give the Ritz bounds.
Every vector in $R$ has native MPS bond at most
$10^6(b_*N)^2(n+1)M^2$. The factor $M^2$ is retained for the whole span.

The formal statement is

\noindent\leanid{TNLean.MPS.AreaLaw.exists_physical_dg_tree_ritz_approximation}.\par
It has no supplied initial viable space, reference, improving family or
trial-energy bound. The tree and final iteration use one common threshold;
an unrelated MPS approximant is not substituted for the root.

This completes the final spectral extraction for the stated modified
construction. It does not identify $j_*$ with the printed iteration
count in terms of $\eta$, prove the interleaved merge procedure or the
energy-estimation invariant, or establish adaptive success probability
or running time. The correction to the printed final-filter gain is
recorded in
\href{https://sirui-lu.com/TNLean/paper-gaps/alvv16_final_filter_gain.pdf}{the final-filter note}.

\section*{A uniform polynomial bound for that same output}

For fixed local dimension $q$ and gap lower bound $\gamma>0$, the
recorded bond of the same modified-tree Ritz output is now bounded by
$CZ^k$, where $C>0$ and $k\in\mathbb N$ precede all physical inputs and
$Z=(L+1)(r+1)(1+\eta^{-1})$. The two actual reference bounds have
profiles at most $b_0+\log(Z^5)$; the inverse refinement parameters have
cubic and quartic bounds in $Z$. With $\beta=1/2$, the actual sample
bound is quadratic in $r$, so the recorded generator count is at most
a constant times $Z^6$. The final natural bond is therefore bounded
by a constant times $Z^{10p_0+23}$, for an integer $p_0$ depending only
on the fixed physical parameters and universal reference constants.

The formal statement is
\leanid{TNLean.MPS.AreaLaw.exists_physical_dg_tree_ritz_polynomial_approximation}.
Its proof extracts the actual operation, balanced root, final literal
iterates and compressed Ritz space from the explicit parent theorem.
It does not replace that output by an independent MPS approximation.
The same space has dimension $r$, mutual error parameter $\eta/2$ and
one MPS bound for its entire span. This is a polynomial bound on the
output witnesses of the stated modified construction; the source
energy-estimation invariant, adaptive probability, source-sharp
parameter dependence and computational running time remain separate.

\section*{The interleaved assignment and integer round count}

In Figure~\texttt{fig:merge-prime}, source lines 1668--1670, the procedure
initializes $V=W'$ and then prints the same assignment
$V=\operatorname{Trim}_{\xi}(\operatorname{span}\{A_iW':i\})$ in every
round. Its right-hand side is independent of the current $V$, so every
positive number of repetitions returns the same subspace. The subsequent
analysis, which compounds the error reduction and the dimension growth
across rounds, requires the intended recursive assignment
\[
  V_0=W',\qquad
  V_{j+1}=\operatorname{Trim}_{\xi}
    \bigl(\operatorname{span}\{A_iV_j:i\}\bigr).
\]
This is a local correction of the assignment, not an identification of a
single trim after an operator-word image with those interleaved trims.
An analysis of the corrected procedure must control the actual trim and
viability at every round. In particular, a constant reference error that
is sufficient after a complete powered-image improvement need not suffice
immediately after a random sample with very small target overlap.

There is a separate integer-round issue in
\texttt{eq:algo-sk-dg}, source line 1797: the printed choice
$k=\tfrac12\lceil\log_Ds\rceil$ need not be an integer, although the
procedure declares $k\in\mathbb N$. For example, $s=D>1$ gives
$k=1/2$. For $D>1$ and $s\ge1$, rounding instead to
$\lceil\tfrac12\log_Ds\rceil$ gives the honest dimension range
$s^2\le D^{2k}s<D^2s^2$, with the larger upper
factor $D^2$. The previously constructed witness uses the explicitly
larger integer power $\lceil\log_D(D^2s)\rceil$ and its own cubic count;
its stated bounds are not the printed count with a missing rounding
argument. Neither correction alone establishes the adaptive probability,
energy-estimation invariant or computational cost of the source algorithm.

\bibliographystyle{unsrt}
\subsection*{The literal interleaved interval tree}
A separate construction now uses the current-space family image followed
by collective trimming at every round of every long parent. Its initial
Haar event is measurable and nonempty, has failure at most $(L+1)^{-2}$
for fixed children, and yields the actual recursive interval tree.
Every retained interval space is $3/200$-viable for the same original
full ground space, has dimension at most $M=D^6s^3$, and has a recorded
spanning family of at most $M$ vectors with the corresponding weighted
internal-cut ranks. The parent equality records an actual chosen Haar
unitary for the actual recursive children.

Here the calibrated window is
\[
 \ell=\left\lceil C\left(v^6/g+\sqrt w\,v^3/\sqrt g\right)\right\rceil,
 \qquad w=1+\log r+2\log(L+1),
\]
with $v=1+\log m+\log(1/g)+\log(1+w)$ and
$D=(m\ell^2)^{22\ell}$. Thus $D$ is honest but length- and rank-dependent.
The construction retains full leaves of length at most $4\ell$, the capped
gap and the explicit integer count $\operatorname{clog}_D(D^2s)$.
This is distinct from the one-final-trim tree described above. The
same interleaved root now has a native matrix product state bound and
literal final refinement. Compressing the original Hamiltonian to that
same refined space gives a Ritz subspace of dimension $r$, both directed
errors at most $\eta/2$, and one bond bound for its entire span. For
fixed $q$ and $\gamma>0$, constants $C>0$ and $k\in\mathbb N$ precede
all physical inputs and give the bound $CZ^k$, where
$Z=(L+1)(r+1)(1+\eta^{-1})$. Moreover, there is a universal
$C_{\rm deg}$ such that
\[
 k\le C_{\rm deg}\frac{(1+\log\max\{q,50\}+\log(1/g))^{10}}{g^2}.
\]
The prefactor can be controlled at the same scale: there is one
universal $C_0>0$, chosen before $q,\gamma$, with
$\log\max\{C,1\}\le C_0(1+\log\max\{q,50\}+\log(1/g))^{10}/g^2$
and the same bound for $k$.
No independent approximation is substituted for the recursive root.
The source energy invariant, adaptive joint probability and running
time remain subsequent conclusions.


\section{The same family for interleaved merging and untrimmed words}

Figure \texttt{fig:algo1}, source lines 1628--1631, uses the family
returned by Generate both for the merge and for the separate untrimmed
energy-estimation image of the child product. A second existence theorem
for an independently chosen family does not identify these operators.

The calibrated physical construction now retains the actual ordered-word
property of its own family, together with its interleaved trimming step.
The finite iteration, fixed-child sampling event and balanced tree retain
these same data. In the tree, the actual target for each long interval
and its quantitative closeness to the original ground space are retained
before every radius and recursive operation. For a $\delta$-viable
input $S$ and every power $n$, the literal ordered-word image has
viability error at most $\Delta^n/(1-\delta)^2$ and dimension at most
$(D^2)^n\dim S$. No second filter is chosen to obtain that conclusion.
The four strengthened statements are
\leanid{TNLean.MPS.AreaLaw.exists_physical_dg_interleaved_filter_family_with_word_powers},
\leanid{TNLean.MPS.AreaLaw.exists_physical_dg_interleaved_interval_iterates_with_word_powers},
\leanid{TNLean.MPS.AreaLaw.exists_physical_dg_interleaved_interval_node_event_with_word_powers}
and
\leanid{TNLean.MPS.AreaLaw.exists_physical_dg_interleaved_interval_tree_with_word_powers}.

The operators in these statements are polynomials in a hard-clipped
Hamiltonian. Their construction is not identified with the soft-truncated
Generate procedure in \texttt{t:aisefficient}. A quantitative transfer
of the actual soft regional energy estimate to this same target can
support a corrected energy candidate; that transfer and the recursive
energy invariant remain separate. The existing sufficient logarithmic
word count is not an identification of the printed
$4\lceil\log(1/\gamma)\rceil$. No running-time claim follows from
retaining the word property. The counterexamples to the printed cluster
and soft-energy estimates are unchanged.

\section{Joint success and the actual canonical Ritz output}

The corrected hard-clipped construction now uses one fixed finite product
of Haar sampling laws for all endpoint pairs and all possible input ranks.
Its true canonical-tree joint success event is Borel. Only the internal
intervals of the actual full-chain subdivision are counted, giving failure
at most $\max\{L-1,0\}/(L+1)^2$; all dimension and recorded-rank bounds
hold independently of success. The physical references, targets, families
and ordered-word certificates are selected together before the common
radius and every seed.

For every requested $\eta>0$, the automatic common radius is chosen before
the seeds. Every successful seed then supplies its own actual root as the
unchanged input to the literal final raw-filter-and-trim iteration. The
compressed Ritz space of that terminal space has dimension $r$, mutual
error $\eta/2$ with the original ground space, and a common all-span bond
$B_{\rm out}\le C((L+1)(r+1)(1+1/\eta))^k$. At fixed $q,\gamma$, the
coefficient and natural degree precede length, multiplicity, accuracy and
seeds. The same bond bound holds for the actual initial native root.

This result retains the hard-clipped family, block leaves, capped-gap
bounds and corrected counts of the preceding sections. It does not
identify that family with the printed soft-truncated Generate procedure,
assert measurability of a chosen Ritz basis, or prove a running-time bound.
The separate energy procedure is not supplied by the Ritz conclusion.

\bibliography{references}
\end{document}
